% Part: first-order-logic
% Chapter: syntax-and-semantics
% Section: first-order-languages

\documentclass[../../../include/open-logic-section]{subfiles}

\begin{document}

\olfileid{fol}{syn}{fol}

\olsection{لغات الرتبة الأولى}


تُبنى تعبيرات منطق الرتبة الأولى من معجم أساسي يحتوي على
\emph{!!{variable}s} و\emph{!!{constant}s} و\emph{!!{predicate}s}
وأحيانًا \emph{!!{function}s}. ومنها، مع الروابط المنطقية والمكممات
وعلامات الترقيم كالأقواس والفواصل، تتكون \emph{الحدود}
و\emph{!!{formula}s}.

\begin{explain}
على نحو غير رسمي، ال!!{predicate}s أسماء للخصائص والعلاقات،
وال!!{constant}s أسماء للكائنات المفردة، وال!!{function}s أسماء
للتطبيقات. وهذه، باستثناء ال!!{identity}~$\eq$، هي
\emph{الرموز غير المنطقية}، وهي تؤلف مجتمعة لغةً ما. وتتحدد أي لغة
رتبة أولى~$\Lang L$ برموزها غير المنطقية. وفي أعم حالة، تحتوي
$\Lang L$ على عدد لا نهائي من الرموز من كل نوع.
\end{explain}

نستعمل في الحالة العامة الرموز الآتية في منطق الرتبة الأولى:

\begin{enumerate}
\item الرموز المنطقية
\begin{enumerate}
\item الروابط المنطقية:
  \startycommalist
  \iftag{prvNot}{\ycomma $\lnot$ (النفي)}{}%
  \iftag{prvAnd}{\ycomma $\land$ (العطف)}{}%
  \iftag{prvOr}{\ycomma $\lor$ (الفصل)}{}%
  \iftag{prvIf}{\ycomma $\lif$ (!!{conditional})}{}%
  \iftag{prvIff}{\ycomma $\liff$ (!!{biconditional})}{}%
  \iftag{prvAll}{\ycomma $\lforall$ (المكمم الكلي)}{}%
  \iftag{prvEx}{\ycomma $\lexists$ (المكمم الوجودي)}{}.
\tagitem{prvFalse}{الثابت القضوي الدال على !!{falsity}~$\lfalse$.}{}
\tagitem{prvTrue}{الثابت القضوي الدال على !!{truth}~$\ltrue$.}{}
\item ال!!{identity} الثنائية المحل~$\eq$.
\item مجموعة !!{denumerable} من ال!!{variable}s: $\Obj v_0$، و$\Obj v_1$، و$\Obj
  v_2$، و\dots
\end{enumerate}
\item الرموز غير المنطقية التي تؤلف \emph{اللغة القياسية} لمنطق
  الرتبة الأولى
\begin{enumerate}
\item مجموعة !!{denumerable} من ال!!{predicate}s ذات $n$ من المحال لكل
  $n>0$: $\Obj A^n_0$، و$\Obj A^n_1$، و$\Obj A^n_2$، و\dots
\item مجموعة !!{denumerable} من ال!!{constant}s: $\Obj c_0$، و$\Obj c_1$، و$\Obj
  c_2$، و\dots.
\item مجموعة !!{denumerable} من ال!!{function}s ذات $n$ من المحال لكل
  $n>0$: $\Obj f^n_0$، و$\Obj f^n_1$، و$\Obj f^n_2$، و\dots
\end{enumerate}
\item علامات الترقيم: القوسان ( و)، والفاصلة.
\end{enumerate}

ستصاغ معظم تعريفاتنا ونتائجنا للغة القياسية الكاملة لمنطق الرتبة
الأولى. غير أننا قد نقصر اللغة، بحسب التطبيق، على عدد قليل من
ال!!{predicate}s وال!!{constant}s وال!!{function}s.

\begin{ex}
تحتوي لغة الحساب~$\Lang L_A$ على !!{predicate} ثنائي المحل واحد~$<$،
و!!{constant} واحد~$\Obj 0$، و!!{function} أحادية المحل واحدة~$\prime$،
و!!{function}s ثنائية المحل اثنتين~$+$ و~$\times$.
\end{ex}

\begin{ex}
لا تحتوي لغة نظرية المجموعات~$\Lang L_Z$ إلا على
ال!!{predicate} الثنائي المحل~$\in$.
\end{ex}

\begin{ex}
لا تحتوي لغة الترتيبات~$\Lang L_\le$ إلا على
ال!!{predicate} الثنائي المحل~$\le$.
\end{ex}

وهذه، مرة أخرى، اصطلاحات: فهي رسميًا مجرد أسماء مستعارة؛ فمثلًا، $<$
و$\in$ و$\le$ أسماء مستعارة لـ~$\Obj A^2_0$، و$\Obj 0$ اسم مستعار
لـ~$\Obj c_0$، و$\prime$ لـ~$\Obj f^1_0$، و$+$ لـ~$\Obj f^2_0$،
و$\times$ لـ~$\Obj f^2_1$.

\iftag{defNot,defOr,defAnd,defIf,defIff,defTrue,defFalse,defEx,defAll}{%
بالإضافة إلى الروابط الأولية و
\iftag{notprvEx,notprvAll}{المكمم}{المكممات} المقدمة أعلاه، نستعمل
أيضًا الرموز \emph{المعرّفة} الآتية:
\startycommalist
  \iftag{defNot}{\ycomma $\lnot$ (النفي)}{}%
  \iftag{defAnd}{\ycomma $\land$ (العطف)}{}%
  \iftag{defOr}{\ycomma $\lor$ (الفصل)}{}%
  \iftag{defIf}{\ycomma $\lif$ (!!{conditional})}{}%
  \iftag{defIff}{\ycomma $\liff$ (!!{biconditional})}{}%
  \iftag{defAll}{\ycomma $\lforall$ (المكمم الكلي)}{}%
  \iftag{defEx}{\ycomma $\lexists$ (المكمم الوجودي)}{}%
  \iftag{defFalse}{\ycomma !!{falsity}~$\lfalse$}{}%
  \iftag{defTrue}{\ycomma !!{truth}~$\ltrue$}{}%
}{}.

\begin{tagblock}{defNot,defOr,defAnd,defIf,defIff,defTrue,defFalse,defEx,defAll}
\begin{explain}
الرمز المعرّف ليس جزءًا رسميًا من اللغة، بل يُدخل بوصفه اختصارًا غير
رسمي: فهو يتيح لنا اختصار صيغ كانت ستغدو طويلة جدًا لو اقتصرنا على
الرموز الأولية. وهذه ميزة واضحة. غير أن الميزة الأكبر هي أن البراهين
تصير أقصر. فإذا كان رمز ما أوليًا، وجب تناوله على حدة في البراهين.
ومن ثم، كلما كثرت الرموز الأولية طالت براهيننا.
\end{explain}
\end{tagblock}

% Alternate symbols

\begin{intro}
قد تكون على معرفة بمصطلحات ورموز تختلف عن تلك التي نستعملها أعلاه. يشيع
في كتب المنطق (وعند مدرّسيه) استعمال $\sim$ و$\neg$ و~!\ للدلالة على
«النفي»، واستعمال $\wedge$ و$\cdot$ و$\&$ للدلالة على «العطف». أما
الرموز الشائعة لـ«الشرط» أو «الاستلزام» فهي $\rightarrow$
و$\Rightarrow$ و$\supset$.
\iftag{prvIff,defIff}{أما رموز «الشرط الثنائي» أو «الاستلزام المتبادل»
  أو «التكافؤ (المادي)» فهي $\leftrightarrow$ و$\Leftrightarrow$
  و$\equiv$.}{}
\iftag{prvFalse,defFalse}{تتنوع أسماء الرمز~$\lfalse$؛ فيسمى
  «الكذب» أو «الباطل» أو «المحال» أو «القاع».}{}
\iftag{prvTrue,defTrue}{تتنوع أسماء الرمز~$\ltrue$؛ فيسمى
  «الصدق» أو «الحق» أو «القمة».}{}

ومن الاصطلاح استعمال الأحرف الصغيرة من أول الأبجدية اللاتينية (مثل
$a$ و$b$ و$c$) لل!!{constant}s (وتسمى أحيانًا أسماء)، والأحرف
الصغيرة من آخرها (مثل $x$ و$y$ و$z$) لل!!{variable}s. وتقترن
المكممات بال!!{variable}s، مثل~$x$؛ ومن الاختلافات الترميزية
$\forall x$ و$(\forall x)$ و$(x)$ و$\Pi x$ و$\bigwedge_x$ للمكمم
الكلي، و$\exists x$ و$(\exists x)$ و$(Ex)$ و$\Sigma x$ و
$\bigvee_x$ للمكمم الوجودي.
\end{intro}

\begin{explain}
قد نعامل جميع المؤثرات القضوية والمكممين كرموز أولية في اللغة. وقد
نختار بدلًا من ذلك ذخيرة أصغر من الرموز الأولية ونعامل سائر
ال!!{operator}s بوصفها معرّفة. ومن مجموعات المؤثرات البوليانية
«الكاملة من حيث الدوال الصدقية» $\{ \lnot, \lor \}$ و
$\{ \lnot, \land \}$ و$\{ \lnot, \lif\}$؛ ويمكن ضم كل منها إلى أي
من المكممين للحصول على لغة رتبة أولى كاملة القدرة التعبيرية.

وقد تكون على معرفة ب!!{operator}s أخرى: شرطة شيفر~$|$ (المنسوبة إلى
هنري شيفر)، وسهم بيرس~$\downarrow$، المعروف أيضًا بخنجر كواين. وإذا
أُعطيا قراءتيهما المعتادتين، «نفي العطف» و«نفي الفصل» (على الترتيب)،
كان كل من هذين المؤثرين كاملًا من حيث الدوال الصدقية بمفرده.
\end{explain}

\end{document}
